\documentclass[12pt]{article}

\usepackage[a4paper, margin=2.5cm]{geometry}
\usepackage{amsmath}
\usepackage{amssymb}
\usepackage{amsthm}
\usepackage[colorlinks=true, linkcolor=blue, citecolor=blue, urlcolor=blue]{hyperref}
\usepackage{booktabs}
\usepackage{natbib}

\newtheorem{theorem}{Theorem}
\newtheorem{proposition}[theorem]{Proposition}
\newtheorem{corollary}[theorem]{Corollary}
\newtheorem{definition}[theorem]{Definition}
\newtheorem{remark}[theorem]{Remark}
\newtheorem{example}[theorem]{Example}

\newcommand{\Mcl}{\mathcal{M}_{\mathrm{cl}}}
\newcommand{\Ccal}{\mathcal{C}}
\newcommand{\Fcal}{\mathcal{F}}
\newcommand{\Hcal}{\mathcal{H}}
\newcommand{\Lcal}{\mathcal{L}}
\newcommand{\Sadm}{\mathcal{S}_{\mathrm{adm}}}
\newcommand{\Scl}{\mathcal{S}_{\mathrm{cl}}}
\newcommand{\phig}{\varphi}

\title{Toward Unitary Closure Dynamics:\\Oscillatory Limits of the Complexified Closure Path Integral}
\author{%
  Lee Smart\\[4pt]
  \textit{Institute of Vibrational Field Dynamics}\\[2pt]
  \texttt{contact@vibrationalfielddynamics.org}\\[2pt]
  \texttt{@vfd\_org}%
}
\date{April 2026}

\begin{document}

\maketitle

\noindent\textbf{Status:} Preprint (not peer-reviewed)

\begin{abstract}
Paper~XV introduced a complexified closure path integral that produces interference from phase-augmented closure dynamics. However, the path amplitude retains a real damping factor alongside the oscillatory phase, so the dynamics is not yet unitary. The present paper derives the differential evolution equation generated by the complexified path integral, identifies its structure as a Fokker--Planck--Schr\"odinger hybrid, and analyses three routes toward an effectively oscillatory regime.

The main result is the \textit{closure evolution equation}: $\partial_t \Psi = \tfrac{\sigma^2}{2}\Delta\Psi - \nabla\Fcal \cdot \nabla\Psi + \tfrac{i}{\sigma^2}\Fcal\,\Psi$, which is derived from the short-time propagator of Paper~XV's complexified path integral. A drift-elimination transformation reveals an operator structure that can be read as Euclidean quantum mechanics plus an imaginary (phase-generating) potential. Three candidate routes to effective oscillatory dynamics are analysed: (A)~measure absorption, (B)~the WKB transport limit $\sigma \to 0$ in which phase dominates and the probability density is transported with local modulus preservation along gradient-flow lines, and (C)~formal Wick rotation. A worked quadratic example demonstrates each regime explicitly.

The paper identifies the precise structural gap between the closure evolution equation and the Schr\"odinger equation: the closure framework can be interpreted as realising the first of two structural complexification steps on one route toward Schr\"odinger form (the imaginary potential from Paper~XV's ansatz) but not the second (imaginary kinetic term). Within the transport regime, local modulus preservation holds along gradient-flow lines, providing a restricted form of effective unitarity; global norm conservation requires additional structure not yet present in the framework.
\end{abstract}

%==================================================================
\section{Introduction}\label{sec:intro}
%==================================================================

Paper~XV~\citep{SmartXV} established the first phase-bearing dynamical layer within the VFD closure programme: a complexified Onsager--Machlup action that produces interference from path-dependent phase accumulation. That paper resolved the missing-phase problem left by Paper~XIV~\citep{SmartXIV} and demonstrated interference as a structural consequence of the complexified path integral.

However, Paper~XV's amplitude retains a real damping factor:
\begin{equation}\label{eq:XV_amplitude}
    \mathcal{A}[\psi(\cdot)] = \exp\!\left(-\frac{S[\psi]}{\sigma^2}\right) \cdot \exp\!\left(\frac{i\, \Fcal_{\mathrm{path}}[\psi]}{\sigma^2}\right)
\end{equation}
where $S[\psi] = \int_0^T \|\dot{\psi} + \nabla\Fcal(\psi)\|^2\, dt$ is the Onsager--Machlup action and $\Fcal_{\mathrm{path}}[\psi] = \int_0^T \Fcal(\psi(t))\, dt$ is the time-integrated closure functional. The first factor is real and dissipative; the second is oscillatory. Standard quantum mechanics is purely oscillatory at the propagator level. The retained damping is the obstacle to unitarity.

\subsection*{The problems this paper addresses}

\begin{enumerate}
    \item[(P5)] \textbf{The retained-damping problem.} Paper~XV introduced phase but retained real damping. Under what conditions does the damping become subdominant, leaving an effectively oscillatory dynamics?

    \item[(P6)] \textbf{The propagator problem.} Paper~XV's path integral defines transition amplitudes. Can a differential evolution equation be derived from the short-time propagator?

    \item[(P7)] \textbf{The effective-unitarity problem.} In quantum mechanics, $|\psi|^2$ is globally conserved. Is there a regime in which the closure framework achieves effective norm conservation?
\end{enumerate}

\subsection*{Scope and epistemic status}

This paper derives the evolution equation generated by the complexified closure path integral and analyses three routes toward oscillatory dynamics. It does not claim full Schr\"odinger recovery, exact unitarity, or Hilbert-space structure. Where results are derived within the adopted formulation, they are stated as such; where routes are formal or exploratory, that status is made explicit.

%==================================================================
\section{Starting Point: Amplitude Decomposition}\label{sec:starting}
%==================================================================

From Paper~XV, the complexified path amplitude is:
\begin{equation}
    \mathcal{A}[\psi(\cdot)] = e^{-S[\psi]/\sigma^2}\, e^{i\Theta[\psi]}
\end{equation}
where the accumulated phase is:
\begin{equation}
    \Theta[\psi(\cdot)] = \frac{1}{\sigma^2}\int_0^T \Fcal(\psi(t))\, dt
\end{equation}

Define the \textit{dissipation-to-phase ratio} along a path:
\begin{equation}\label{eq:epsilon}
    \varepsilon[\psi] = \frac{S[\psi]}{\Fcal_{\mathrm{path}}[\psi]} = \frac{\int_0^T \|\dot{\psi} + \nabla\Fcal\|^2\, dt}{\int_0^T \Fcal(\psi)\, dt}
\end{equation}

This ratio characterises the dynamical regime:
\begin{itemize}
    \item $\varepsilon \gg 1$: dissipation-dominated (Paper~XIV regime),
    \item $\varepsilon \sim 1$: mixed regime (Paper~XV interference),
    \item $\varepsilon \ll 1$: phase-dominated (oscillatory target of this paper).
\end{itemize}

The transition propagator is:
\begin{equation}\label{eq:propagator}
    K(\psi_f, \psi_i; T) = \int_{\psi(0)=\psi_i}^{\psi(T)=\psi_f} \mathcal{A}[\psi(\cdot)]\, \mathcal{D}\psi
\end{equation}

and the state evolves as:
\begin{equation}\label{eq:evolution}
    \Psi(\psi, t+\Delta t) = \int K(\psi, \psi'; \Delta t)\, \Psi(\psi', t)\, d\psi'
\end{equation}

%==================================================================
\section{The Closure Evolution Equation}\label{sec:evolution}
%==================================================================

The central result of this paper is the differential equation that the propagator~\eqref{eq:propagator} generates for the evolving state $\Psi$.

\subsection{Derivation from the short-time propagator}

The stochastic dynamics of Paper~XIV is governed by the Langevin equation:
\begin{equation}
    d\psi_t = -\nabla\Fcal(\psi_t)\, dt + \sigma\, dW_t
\end{equation}
whose backward Kolmogorov generator is:
\begin{equation}
    L_{\mathrm{BK}} = \frac{\sigma^2}{2}\,\Delta - \nabla\Fcal \cdot \nabla
\end{equation}

This generator produces the short-time transition kernel $G(\psi, \psi'; \Delta t)$ for the \textit{real} (phase-free) dynamics. Paper~XV's complexification adds a phase factor $\exp(i\Fcal(\psi')\Delta t/\sigma^2)$ per time step, evaluated at the starting configuration. The complexified kernel is therefore:
\begin{equation}
    K(\psi, \psi'; \Delta t) = G(\psi, \psi'; \Delta t) \cdot \exp\!\left(\frac{i\, \Fcal(\psi')\, \Delta t}{\sigma^2}\right)
\end{equation}

Expanding to first order in $\Delta t$:
\begin{align}
    \Psi(\psi, t+\Delta t) &= \int G(\psi, \psi'; \Delta t)\, e^{i\Fcal(\psi')\Delta t/\sigma^2}\, \Psi(\psi', t)\, d\psi' \notag \\
    &\approx (1 + \Delta t\, L_{\mathrm{BK}})\!\left[\left(1 + \frac{i\Fcal\,\Delta t}{\sigma^2}\right)\Psi\right] \notag \\
    &= \Psi + \Delta t\left[L_{\mathrm{BK}}\Psi + \frac{i\Fcal}{\sigma^2}\,\Psi\right] + O(\Delta t^2)
\end{align}

This yields the \textit{closure evolution equation}:

\begin{proposition}[Closure evolution equation]\label{prop:evolution}
Under the formal short-time-generator expansion and sufficient regularity of $\Fcal$ (Remark~\ref{rem:derivation_status}), the complexified closure path integral of Paper~XV generates the differential evolution equation:
\begin{equation}\label{eq:closure_evolution}
    \boxed{\partial_t \Psi = \frac{\sigma^2}{2}\,\Delta\Psi - \nabla\Fcal \cdot \nabla\Psi + \frac{i}{\sigma^2}\,\Fcal\,\Psi}
\end{equation}
equivalently:
\begin{equation}
    \partial_t \Psi = L_{\mathrm{BK}}\Psi + \frac{i}{\sigma^2}\,\Fcal\,\Psi
\end{equation}
\end{proposition}

\begin{remark}[Derivation status]\label{rem:derivation_status}
The derivation is formal: it relies on the standard short-time kernel expansion of the backward Kolmogorov semigroup generated by the Paper~XIV Langevin equation, together with sufficient smoothness of $\Fcal$ (at minimum $C^2$ with bounded derivatives on the relevant domain). A fully rigorous construction of the associated complexified semigroup --- including domain specification, strong continuity, and well-posedness of the initial-value problem for~\eqref{eq:closure_evolution} --- is not attempted here and is left for future work. The consistency checks below and the explicit quadratic solution in Section~\ref{sec:example} provide supporting evidence for the formal result.
\end{remark}

The three terms have distinct physical roles:
\begin{enumerate}
    \item $\frac{\sigma^2}{2}\Delta\Psi$: \textbf{diffusion} from the stochastic noise of Paper~XIV.
    \item $-\nabla\Fcal \cdot \nabla\Psi$: \textbf{drift} toward closure (gradient-flow transport from Paper~XII).
    \item $\frac{i}{\sigma^2}\Fcal\,\Psi$: \textbf{phase rotation} proportional to the closure functional (from Paper~XV's complexification ansatz).
\end{enumerate}

\begin{remark}[Structural identification]
The closure evolution equation~\eqref{eq:closure_evolution} is a \textit{Fokker--Planck--Schr\"odinger hybrid}: the first two terms form the backward Kolmogorov generator of the stochastic closure dynamics (Papers~XII--XIV), and the third term is the imaginary potential introduced by Paper~XV's complexification ansatz.
\end{remark}

\subsection{Consistency checks}

\textbf{Phase-free limit ($\Fcal_{\mathrm{path}} \to 0$).} Setting the imaginary potential to zero recovers:
\begin{equation}
    \partial_t \Psi = \frac{\sigma^2}{2}\Delta\Psi - \nabla\Fcal \cdot \nabla\Psi = L_{\mathrm{BK}}\Psi
\end{equation}
which is the backward Kolmogorov equation of Paper~XIV. $\checkmark$

\textbf{Zero-noise limit ($\sigma \to 0$).} The diffusion term vanishes, leaving:
\begin{equation}
    \partial_t \Psi = -\nabla\Fcal \cdot \nabla\Psi + \frac{i}{\sigma^2}\Fcal\,\Psi
\end{equation}
which is deterministic gradient-flow transport with rapid phase rotation --- consistent with Paper~XII dynamics plus the phase structure of Paper~XV. $\checkmark$

\textbf{On the closure manifold ($\Fcal = 0$).} At any point where $\Fcal(\psi) = 0$ and $\nabla\Fcal(\psi) = 0$:
\begin{equation}
    \partial_t \Psi = \frac{\sigma^2}{2}\Delta\Psi
\end{equation}
Pure diffusion with no phase rotation and no drift. Closed states diffuse but do not rotate. $\checkmark$

%==================================================================
\section{Drift Elimination and Structural Analysis}\label{sec:drift}
%==================================================================

The drift term $-\nabla\Fcal \cdot \nabla\Psi$ can be removed by a change of variables, revealing the underlying operator structure.

\begin{proposition}[Drift-free form]\label{prop:driftfree}
Under the substitution $\Psi(\psi,t) = e^{\Fcal(\psi)/\sigma^2}\, \Phi(\psi,t)$, the closure evolution equation~\eqref{eq:closure_evolution} becomes:
\begin{equation}\label{eq:driftfree}
    \partial_t \Phi = \frac{\sigma^2}{2}\,\Delta\Phi + W(\psi)\,\Phi
\end{equation}
where the complex effective potential is:
\begin{equation}\label{eq:W}
    W(\psi) = -\frac{|\nabla\Fcal|^2}{2\sigma^2} + \frac{\Delta\Fcal}{2} + \frac{i\,\Fcal}{\sigma^2}
\end{equation}
\end{proposition}

\begin{proof}
Substitute $\Psi = e^{\Fcal/\sigma^2}\Phi$ into~\eqref{eq:closure_evolution}. The left-hand side gives $e^{\Fcal/\sigma^2}\partial_t\Phi$. For the Laplacian:
\begin{equation}
    \Delta(e^{\Fcal/\sigma^2}\Phi) = e^{\Fcal/\sigma^2}\!\left[\Delta\Phi + \frac{2}{\sigma^2}\nabla\Fcal\cdot\nabla\Phi + \frac{|\nabla\Fcal|^2}{\sigma^4}\Phi + \frac{\Delta\Fcal}{\sigma^2}\Phi\right]
\end{equation}
For the drift:
\begin{equation}
    -\nabla\Fcal\cdot\nabla(e^{\Fcal/\sigma^2}\Phi) = -e^{\Fcal/\sigma^2}\!\left[\frac{|\nabla\Fcal|^2}{\sigma^2}\Phi + \nabla\Fcal\cdot\nabla\Phi\right]
\end{equation}
Collecting all terms and dividing by $e^{\Fcal/\sigma^2}$:
\begin{align}
    \partial_t\Phi &= \frac{\sigma^2}{2}\Delta\Phi + \nabla\Fcal\cdot\nabla\Phi + \frac{|\nabla\Fcal|^2}{2\sigma^2}\Phi + \frac{\Delta\Fcal}{2}\Phi \notag \\
    &\quad - \frac{|\nabla\Fcal|^2}{\sigma^2}\Phi - \nabla\Fcal\cdot\nabla\Phi + \frac{i\Fcal}{\sigma^2}\Phi
\end{align}
The $\nabla\Fcal\cdot\nabla\Phi$ terms cancel, yielding~\eqref{eq:driftfree} with $W$ as in~\eqref{eq:W}.
\end{proof}

\begin{remark}[Witten--supersymmetric connection]\label{rem:witten}
The real part of $W$ is:
\begin{equation}
    \mathrm{Re}(W) = -\frac{|\nabla\Fcal|^2}{2\sigma^2} + \frac{\Delta\Fcal}{2} = -V_{\mathrm{W}}(\psi)
\end{equation}
where $V_{\mathrm{W}} = |\nabla\Fcal|^2/(2\sigma^2) - \Delta\Fcal/2$ is the \textit{Witten potential} familiar from supersymmetric quantum mechanics. In the Witten framework, the ground-state wavefunction of the supersymmetric partner Hamiltonian is $\psi_0 \propto e^{-\Fcal/\sigma^2}$ --- precisely the Paper~XIV equilibrium amplitude. The drift-free form~\eqref{eq:driftfree} therefore embeds the closure dynamics within a structure that already has known connections to quantum mechanics, though those connections are not invoked here as physical identifications. The relevance of the Witten potential is operator-theoretic rather than identificatory: it clarifies the transformed real part of the closure generator without implying that the closure framework is itself a supersymmetric quantum system.
\end{remark}

The drift-free equation~\eqref{eq:driftfree} separates cleanly into:
\begin{equation}
    \partial_t \Phi = \underbrace{\frac{\sigma^2}{2}\Delta\Phi - V_{\mathrm{W}}\Phi}_{\text{Euclidean quantum mechanics}} + \underbrace{\frac{i\Fcal}{\sigma^2}\Phi}_{\text{phase rotation (from XV)}}
\end{equation}

This decomposition reveals the operator structure of the closure evolution: it has the operator form of \textbf{Euclidean quantum mechanics plus an imaginary potential} introduced by Paper~XV's complexification ansatz.

%==================================================================
\section{Routes to Effective Oscillatory Dynamics}\label{sec:routes}
%==================================================================

Three routes toward an effectively oscillatory regime are analysed.

\subsection{Route A: Measure absorption}\label{sec:routeA}

The path integral~\eqref{eq:propagator} can be formally rewritten:
\begin{equation}
    K(\psi_f, \psi_i; T) = \int e^{i\Theta[\psi]/\sigma^2}\, d\mu[\psi]
\end{equation}
where:
\begin{equation}
    d\mu[\psi] = e^{-S[\psi]/\sigma^2}\, \mathcal{D}\psi
\end{equation}
is the \textit{Onsager--Machlup measure} --- a positive, real-valued weighted path measure.

In this rewriting, the propagator is a purely oscillatory integral over a weighted measure. The damping is not eliminated; it is absorbed into the definition of which paths contribute significantly.

\begin{remark}
Measure absorption is a formal rearrangement, not a dynamical simplification. The measure $d\mu$ is concentrated on paths close to the gradient-flow solution (those with small Onsager--Machlup action), and the phase $\Theta$ then generates oscillatory contributions among those dynamically favoured paths. This is the natural reading of the closure path integral: \textit{the closure structure selects which paths matter; the phase determines how they interfere.}
\end{remark}

\subsection{Route B: WKB transport regime ($\sigma \to 0$)}\label{sec:routeB}

In the limit $\sigma \to 0$, the diffusion term in~\eqref{eq:closure_evolution} vanishes while the phase term $i\Fcal/\sigma^2$ diverges. The dominant dynamics becomes:

\begin{equation}\label{eq:transport}
    \partial_t \Psi = -\nabla\Fcal \cdot \nabla\Psi + \frac{i}{\sigma^2}\,\Fcal\,\Psi
\end{equation}

This is a \textit{transport equation with phase rotation}: the configuration is advected along the gradient-flow lines $\dot{\psi} = -\nabla\Fcal$ (Paper~XII), and the state acquires phase at rate $\Fcal/\sigma^2$. This regime is obtained by neglecting the diffusive term relative to the drift and phase terms in the asymptotic limit $\sigma \to 0$; the results that follow are valid within this transport approximation.

\begin{proposition}[Norm transport in the WKB regime]\label{prop:normtransport}
In the transport regime~\eqref{eq:transport}, the probability density $\rho = |\Psi|^2$ satisfies:
\begin{equation}\label{eq:rhotransport}
    \partial_t \rho + \nabla\Fcal \cdot \nabla\rho = 0
\end{equation}
That is, $\rho$ is conserved along gradient-flow lines: $\frac{D\rho}{Dt} = 0$ where $\frac{D}{Dt} = \partial_t - \nabla\Fcal \cdot \nabla$ is the material derivative along $\dot{\psi} = -\nabla\Fcal$.
\end{proposition}

\begin{proof}
From~\eqref{eq:transport}:
\begin{align}
    \partial_t|\Psi|^2 &= \Psi^*\partial_t\Psi + \Psi\,\partial_t\Psi^* \notag \\
    &= \Psi^*\!\left(-\nabla\Fcal\cdot\nabla\Psi + \frac{i\Fcal}{\sigma^2}\Psi\right) + \Psi\!\left(-\nabla\Fcal\cdot\nabla\Psi^* - \frac{i\Fcal}{\sigma^2}\Psi^*\right) \notag \\
    &= -\nabla\Fcal \cdot \nabla|\Psi|^2
\end{align}
where the phase terms cancel: $\Psi^*(i\Fcal/\sigma^2)\Psi + \Psi(-i\Fcal/\sigma^2)\Psi^* = 0$.
\end{proof}

This is the key modulus-preservation result of the paper: \textbf{in the WKB transport approximation, the phase rotation does not locally alter the probability density.} The modulus $|\Psi|$ is advected along gradient-flow characteristics without phase-induced gain or loss. This is \textit{local} conservation along individual flow lines, not global unitarity in the Schr\"odinger sense --- the distinction is discussed in Remark~\ref{rem:global}.

\begin{remark}[Global vs.\ local conservation]\label{rem:global}
Proposition~\ref{prop:normtransport} establishes \textit{local} conservation: $D\rho/Dt = 0$ along each flow line. However, the total norm $\int |\Psi|^2\, d\psi$ satisfies:
\begin{equation}
    \partial_t \int |\Psi|^2\, d\psi = \int \Delta\Fcal\, |\Psi|^2\, d\psi
\end{equation}
which vanishes only if $\Delta\Fcal = 0$ everywhere (harmonic $\Fcal$) or if $|\Psi|^2$ is supported where $\Delta\Fcal = 0$. The gradient flow $\dot\psi = -\nabla\Fcal$ is generally compressible ($\nabla \cdot (-\nabla\Fcal) = -\Delta\Fcal \neq 0$), so global norm conservation is \textit{not} automatic. This is a structural difference from the Schr\"odinger equation, where unitarity gives global norm conservation exactly.
\end{remark}

\subsection{Route C: Formal Wick rotation}\label{sec:routeC}

The drift-free equation~\eqref{eq:driftfree} has the structure of a real diffusion equation with a complex potential. A formal Wick rotation $t \to -i\tau$ converts the real-time diffusion to an imaginary-time ``Schr\"odinger-like'' form.

Under $t = -i\tau$ (so $\partial_t = i\,\partial_\tau$):
\begin{equation}\label{eq:wick}
    i\,\partial_\tau \Phi = \frac{\sigma^2}{2}\,\Delta\Phi + W(\psi)\,\Phi
\end{equation}

Comparing with the Schr\"odinger equation $i\hbar\,\partial_t\psi = -(\hbar^2/2m)\Delta\psi + V\psi$:
\begin{itemize}
    \item the kinetic term has $+\sigma^2/2$ (positive), whereas Schr\"odinger has $-\hbar^2/(2m)$ (negative),
    \item the potential is complex: $W = -V_{\mathrm{W}} + i\Fcal/\sigma^2$.
\end{itemize}

The sign of the kinetic term is opposite to the Schr\"odinger convention. In the near-closure regime where $V_{\mathrm{W}} \approx 0$, equation~\eqref{eq:wick} becomes:
\begin{equation}
    i\,\partial_\tau \Phi \approx \frac{\sigma^2}{2}\,\Delta\Phi + \frac{i\Fcal}{\sigma^2}\,\Phi
\end{equation}

A further formal substitution $\sigma^2 \to -\sigma^2$ (or equivalently $\sigma \to i\tilde{\sigma}$) would flip the kinetic sign:
\begin{equation}
    i\,\partial_\tau \Phi \to -\frac{\tilde{\sigma}^2}{2}\,\Delta\Phi + \frac{i\Fcal}{i\tilde{\sigma}^2}\Phi = -\frac{\tilde{\sigma}^2}{2}\,\Delta\Phi + \frac{\Fcal}{\tilde{\sigma}^2}\,\Phi
\end{equation}

Under this formal double analytic continuation, the result has the form of the Schr\"odinger equation with $\hbar = 1$, $m = 1/\tilde{\sigma}^2$, and potential $V = \Fcal/\tilde{\sigma}^2$. This is a formal algebraic observation; neither substitution is physically derived from the closure framework.

\begin{remark}[Status of the Wick route]
The formal double substitution ($t \to -i\tau$, $\sigma^2 \to -\tilde{\sigma}^2$) maps the closure evolution equation to the Schr\"odinger equation. This is a structural observation, not a derived physical correspondence. Two formal operations are required: the time rotation and the noise-parameter continuation. Neither is justified by the closure framework alone; they are noted as a possible route to exact Schr\"odinger recovery in future work. The physically more robust result is Route~B (WKB transport), which requires no analytic continuation.
\end{remark}

\subsection{Summary: the two-complexification structure}

Paper~XV performed one complexification: adding an imaginary potential $i\Fcal/\sigma^2$ to the path action. This produced interference.

The present analysis suggests a structural two-step complexification route: a second formal complexification --- of the kinetic term, via $\sigma^2 \to i\tilde{\sigma}^2$ or equivalently a time rotation --- would formally produce the Schr\"odinger equation. The two steps are:

\begin{table}[ht]
\centering
\caption{Structural two-step complexification route from closure dynamics toward Schr\"odinger form.}
\label{tab:twostep}
\begin{tabular}{@{}lll@{}}
\toprule
Step & What is complexified & What is gained \\
\midrule
Paper~XV & Potential: $S \to S - i\Fcal_{\mathrm{path}}$ & Phase, interference \\
Paper~XVI (formal) & Kinetic: $\sigma^2 \to -\tilde{\sigma}^2$ & Schr\"odinger structure \\
\bottomrule
\end{tabular}
\end{table}

The first step is the content of Paper~XV and is adopted throughout this paper. The second step is formal and unexplored; it is recorded here as a structural observation.

%==================================================================
\section{Effective Unitarity}\label{sec:unitarity}
%==================================================================

In quantum mechanics, unitarity means $\partial_t \int |\psi|^2 = 0$ globally. The closure framework does not achieve this exactly. We introduce a working criterion for approximate norm conservation and characterise the regimes in which it holds.

\textbf{Working criterion (effective unitarity).}\label{def:effunitary} For the purposes of this paper, we say the closure evolution is \textit{effectively unitary} on a domain $\mathcal{D} \subset \Sadm$ over a timescale $T$ if the norm variation is negligible relative to the dynamical evolution:
\begin{equation}
    \left|\int_0^T \partial_t \int_{\mathcal{D}} |\Psi|^2\, d\psi\, dt\right| \ll \int_{\mathcal{D}} |\Psi(\cdot, 0)|^2\, d\psi
\end{equation}
This is an operational notion, not a formal definition of a unitary group. It identifies regimes where the norm is approximately preserved over the timescale on which the phase dynamics significantly evolves.

From the full evolution equation~\eqref{eq:closure_evolution}, assuming sufficient decay of $\Psi$ at infinity (or compact support) so that boundary terms from integration by parts vanish, the norm evolves as:
\begin{align}\label{eq:normevolution}
    \partial_t \int |\Psi|^2\, d\psi &= \sigma^2 \int \mathrm{Re}(\Psi^*\Delta\Psi)\, d\psi - \int \nabla\Fcal \cdot \nabla|\Psi|^2\, d\psi \notag \\
    &= -\sigma^2 \int |\nabla\Psi|^2\, d\psi + \int \Delta\Fcal\, |\Psi|^2\, d\psi
\end{align} The phase term contributes zero (as shown in Proposition~\ref{prop:normtransport}).

\subsection*{Regimes supporting effective unitarity}

Three regimes in which the total norm $\int|\Psi|^2$ is approximately conserved can be identified from~\eqref{eq:normevolution}:

\textbf{(i) WKB regime} ($\sigma \to 0$). The diffusive-loss term $\sigma^2 \int |\nabla\Psi|^2$ vanishes. Norm variation is then controlled by $\int \Delta\Fcal\, |\Psi|^2$, which is small if $|\Psi|^2$ is concentrated where $\Delta\Fcal \approx 0$. This is the asymptotic regime of Proposition~\ref{prop:normtransport}.

\textbf{(ii) Harmonic closure} ($\Delta\Fcal = 0$). The drift is divergence-free, so $|\Psi|^2$ is exactly transported. The remaining norm evolution is $\partial_t \int |\Psi|^2 = -\sigma^2 \int |\nabla\Psi|^2 \leq 0$: diffusive loss only, with no compressive gain. This is an exact special case.

\textbf{(iii) Cancellation regime.} For specific states $\Psi$, the diffusive loss $\sigma^2 \int |\nabla\Psi|^2$ may be approximately balanced by the compressive gain $\int \Delta\Fcal\, |\Psi|^2$. The Paper~XIV stationary state is the exact fixed point of this balance (confirmed in the quadratic example, Section~\ref{sec:example}). This is a state-dependent, contingent condition, not a structural guarantee.

\begin{remark}
Exact global unitarity ($\partial_t \int |\Psi|^2 = 0$ for all $\Psi$) would require both $\sigma = 0$ and $\Delta\Fcal = 0$ simultaneously --- a condition too restrictive to hold generically. The closure framework achieves effective unitarity in restricted regimes, not exact unitarity as a structural property.
\end{remark}

%==================================================================
\section{Worked Example: Quadratic Closure Functional}\label{sec:example}
%==================================================================

\begin{example}[Quadratic closure functional]\label{ex:quadratic}
Consider the one-dimensional closure functional:
\begin{equation}
    \Fcal(\psi) = \frac{1}{2}k\psi^2, \qquad k > 0
\end{equation}
with $\nabla\Fcal = k\psi$, $\Delta\Fcal = k$, $|\nabla\Fcal|^2 = k^2\psi^2$.

\subsection*{Full evolution equation}

Equation~\eqref{eq:closure_evolution} becomes:
\begin{equation}\label{eq:quadratic_evo}
    \partial_t \Psi = \frac{\sigma^2}{2}\,\partial_\psi^2 \Psi - k\psi\,\partial_\psi\Psi + \frac{ik\psi^2}{2\sigma^2}\,\Psi
\end{equation}

The three terms: diffusion ($\sigma^2/2\cdot\partial^2$), Ornstein--Uhlenbeck drift ($-k\psi\cdot\partial$), and quadratic phase rotation ($ik\psi^2/(2\sigma^2)$).

\subsection*{Drift-free form}

Under $\Psi = e^{k\psi^2/(2\sigma^2)}\Phi$, equation~\eqref{eq:driftfree} gives:
\begin{equation}
    \partial_t \Phi = \frac{\sigma^2}{2}\,\partial_\psi^2 \Phi + \left(-\frac{k^2\psi^2}{2\sigma^2} + \frac{k}{2} + \frac{ik\psi^2}{2\sigma^2}\right)\Phi
\end{equation}

The effective potential is:
\begin{equation}
    W(\psi) = \frac{k}{2} + \frac{k(i - k)}{2\sigma^2}\,\psi^2
\end{equation}

The Witten potential is $V_{\mathrm{W}} = k^2\psi^2/(2\sigma^2) - k/2$.

\subsection*{WKB transport regime}

For $\sigma \to 0$, equation~\eqref{eq:transport} gives:
\begin{equation}
    \partial_t \Psi = -k\psi\,\partial_\psi\Psi + \frac{ik\psi^2}{2\sigma^2}\,\Psi
\end{equation}

Along the characteristic $\psi(t) = \psi_0\, e^{-kt}$:
\begin{equation}
    \Psi(\psi(t), t) = \exp\!\left(\frac{ik}{2\sigma^2}\int_0^t \psi_0^2 e^{-2ks}\, ds\right)\Psi(\psi_0, 0)
    = \exp\!\left(\frac{i\psi_0^2(1-e^{-2kt})}{4\sigma^2}\right)\Psi(\psi_0, 0)
\end{equation}

The modulus $|\Psi|$ is constant along each characteristic: $|\Psi(\psi(t),t)| = |\Psi(\psi_0, 0)|$. The phase accumulates at a rate proportional to the closure cost $\Fcal = k\psi^2/2$, which decays as the configuration approaches closure ($\psi \to 0$).

\subsection*{Norm evolution}

From~\eqref{eq:normevolution}:
\begin{equation}
    \partial_t \int |\Psi|^2\, d\psi = -\sigma^2 \int |\partial_\psi\Psi|^2\, d\psi + k\int |\Psi|^2\, d\psi
\end{equation}

The compressive term $+k\int|\Psi|^2$ (from $\Delta\Fcal = k > 0$) tends to increase the norm (the Ornstein--Uhlenbeck drift concentrates probability). The diffusive term $-\sigma^2\int|\partial_\psi\Psi|^2$ tends to decrease it. For the stationary state $\Psi_{\mathrm{st}} \propto e^{-k\psi^2/(2\sigma^2)}$, these balance: $\sigma^2 \int |\partial_\psi\Psi_{\mathrm{st}}|^2 = k\int|\Psi_{\mathrm{st}}|^2$. This confirms that the Paper~XIV equilibrium is the fixed point of the real part of the evolution.

\subsection*{Regime diagram}

\begin{table}[ht]
\centering
\caption{Regimes for the quadratic example $\Fcal = \frac{1}{2}k\psi^2$.}
\label{tab:quadregimes}
\begin{tabular}{@{}lll@{}}
\toprule
Regime & Dominant terms & Behaviour \\
\midrule
$\sigma^2 \gg k\psi^2$ & Diffusion & Spreading, decoherence \\
$\sigma^2 \sim k\psi^2$ & All three & Mixed: diffusion + drift + phase \\
$\sigma^2 \ll k\psi^2$ & Drift + phase & WKB transport with oscillation \\
\bottomrule
\end{tabular}
\end{table}

\end{example}

%==================================================================
\section{Comparison with the Schr\"odinger Equation}\label{sec:schrodinger}
%==================================================================

The closure evolution equation and the Schr\"odinger equation can now be compared term by term:

\begin{table}[ht]
\centering
\caption{Term-by-term comparison of the closure and Schr\"odinger evolution.}
\label{tab:comparison}
\begin{tabular}{@{}lll@{}}
\toprule
Term & Schr\"odinger ($i\hbar\,\partial_t\psi = \hat{H}\psi$) & Closure ($\partial_t\Psi = \Lcal\Psi$) \\
\midrule
Time derivative & $i\hbar\,\partial_t$ (purely imaginary) & $\partial_t$ (real) \\
Kinetic & $-\frac{\hbar^2}{2m}\Delta$ (real, negative) & $\frac{\sigma^2}{2}\Delta$ (real, positive) \\
Drift & Absent & $-\nabla\Fcal\cdot\nabla$ \\
Potential & $V(\psi)$ (real) & $\frac{i}{\sigma^2}\Fcal$ (imaginary) \\
\bottomrule
\end{tabular}
\end{table}

The structural differences are:
\begin{enumerate}
    \item \textbf{Kinetic sign.} Schr\"odinger has $-\hbar^2/(2m)\Delta$ (after dividing by $i\hbar$, this becomes $+i\hbar/(2m)\Delta$: imaginary diffusion). The closure equation has $+\sigma^2/2\cdot\Delta$: real diffusion. This is the signature of a dissipative rather than unitary kinetic term.

    \item \textbf{Drift term.} The closure equation contains $-\nabla\Fcal\cdot\nabla\Psi$, which has no counterpart in the free Schr\"odinger equation. (It can be eliminated by the substitution of Proposition~\ref{prop:driftfree}, at the cost of introducing the Witten potential.)

    \item \textbf{Potential type.} The Schr\"odinger potential is real; the closure ``potential'' $i\Fcal/\sigma^2$ is imaginary. An imaginary potential in quantum mechanics represents absorption or gain --- here it generates phase rotation.
\end{enumerate}

\begin{remark}[The two-complexification gap]
Sections~\ref{sec:routes} and~\ref{sec:schrodinger} reveal the precise structural gap between the closure evolution and the Schr\"odinger equation. Paper~XV's complexification ansatz performed the first of two structural steps identified on the route toward Schr\"odinger form: making the \textit{potential} imaginary. The second step --- making the \textit{kinetic term} imaginary (i.e.\ converting real diffusion to the imaginary diffusion of quantum mechanics) --- has not been performed and is not justified by the closure framework alone. Route~C (Section~\ref{sec:routeC}) shows that a formal double analytic continuation achieves this algebraically, but the physical interpretation of such a continuation remains open. Whether this two-step route is the only path from closure to Schr\"odinger dynamics, or whether alternative routes exist, is not determined by the present analysis.
\end{remark}

%==================================================================
\section{Discussion and Limitations}\label{sec:discussion}
%==================================================================

\subsection{What is achieved}

\textbf{Entry assumptions.} This paper adopts Paper~XV's complexified-action ansatz and the standard Langevin/Fokker--Planck framework of Paper~XIV.

\textbf{Derived within the adopted formulation:}
\begin{itemize}
    \item The closure evolution equation~\eqref{eq:closure_evolution} is derived in the formal short-time-generator sense from the complexified path integral (Proposition~\ref{prop:evolution}, Remark~\ref{rem:derivation_status}).
    \item The drift-free form~\eqref{eq:driftfree} is derived, revealing an operator structure that can be read as Euclidean QM plus an imaginary potential (Proposition~\ref{prop:driftfree}).
    \item In the WKB transport regime ($\sigma \to 0$), the probability density $|\Psi|^2$ is locally preserved along gradient-flow characteristics (Proposition~\ref{prop:normtransport}). Global norm conservation is not automatic (Remark~\ref{rem:global}).
    \item The norm evolution equation~\eqref{eq:normevolution} identifies the two competing terms: diffusive loss and compressive gain.
    \item The quadratic example is analysed explicitly across the three regimes.
\end{itemize}

\textbf{Structural observations (not derived as physical law):}
\begin{itemize}
    \item The connection to the Witten potential and supersymmetric quantum mechanics (Remark~\ref{rem:witten}).
    \item The structural two-step complexification route toward Schr\"odinger form (Table~\ref{tab:twostep}).
    \item The formal Wick rotation to Schr\"odinger form (Section~\ref{sec:routeC}).
\end{itemize}

\subsection{What is not achieved}

\begin{itemize}
    \item \textbf{No exact unitarity.} Global norm conservation does not hold in general. It requires either $\sigma = 0$ and $\Delta\Fcal = 0$, or a cancellation that is state-dependent.
    \item \textbf{No Schr\"odinger equation.} The closure evolution equation is structurally different from Schr\"odinger: real diffusion vs.\ imaginary diffusion. A formal double analytic continuation maps one to the other, but this is not derived.
    \item \textbf{No Hilbert-space structure.} The augmented state space remains $\Sadm \times S^1$, not a Hilbert space.
    \item \textbf{No physical interpretation of the Wick route.} The formal substitutions $t \to -i\tau$, $\sigma^2 \to -\tilde{\sigma}^2$ are not justified by the closure framework.
    \item \textbf{The quadratic example is schematic.} It demonstrates the mechanism but does not compute predictions for a physical system.
\end{itemize}

\subsection{Open questions}

\begin{itemize}
    \item Is there a physical principle within the closure framework that selects or motivates the second complexification (imaginary kinetic term)?
    \item Can the compressive gain $\int \Delta\Fcal\,|\Psi|^2$ be cancelled by a suitable choice of inner product or norm (e.g.\ a $P_{\mathrm{st}}$-weighted norm)?
    \item Does the Witten-potential connection provide a path to deriving the closure evolution from a supersymmetric Hamiltonian?
    \item Can the WKB transport regime be made precise enough to serve as an effective quantum mechanics for the closure framework?
\end{itemize}

%==================================================================
\section{Conclusion}\label{sec:conclusion}
%==================================================================

This paper solves the propagator problem left by Paper~XV. The complexified closure path integral generates a differential evolution equation (Proposition~\ref{prop:evolution}), derived in the formal short-time-generator sense under standard regularity assumptions. The closure evolution equation is a Fokker--Planck--Schr\"odinger hybrid --- the backward Kolmogorov generator of the stochastic closure dynamics plus the imaginary potential from Paper~XV's complexification ansatz.

Three routes toward effective oscillatory dynamics have been analysed:
\begin{itemize}
    \item \textbf{Measure absorption} reinterprets the damping as a path-measure weight, leaving a formally oscillatory integral over a weighted measure.
    \item \textbf{WKB transport} ($\sigma \to 0$) produces a regime where the probability density is advected along gradient-flow lines with local modulus preservation (Proposition~\ref{prop:normtransport}). This is the physically most robust route: it requires no analytic continuation and holds as an asymptotic result within the adopted framework.
    \item \textbf{Formal Wick rotation} reveals that a formal double analytic continuation maps the closure evolution to the Schr\"odinger equation, identifying a structural two-step complexification route. This is an algebraic observation, not a derived physical correspondence.
\end{itemize}

The honest assessment is as follows. This paper derives the closure evolution equation, identifies its structure, and characterises a WKB regime with local modulus preservation along flow lines. It does \textit{not} derive exact global unitarity (which requires conditions too restrictive to hold generically), does \textit{not} derive the Schr\"odinger equation (the kinetic term retains the wrong sign), and does \textit{not} physically justify the formal Wick route. The structural two-step complexification is identified but not completed.

The narrative arc is now:
\begin{itemize}
    \item Paper~XII: deterministic gradient flow,
    \item Paper~XIII: multi-state interaction,
    \item Paper~XIV: stochastic dynamics and Boltzmann equilibrium,
    \item Paper~XV: phase, interference, squared-modulus probability,
    \item Paper~XVI: evolution equation, structural identification, local modulus preservation in the WKB regime.
\end{itemize}

The next frontier is whether the second complexification step (imaginary kinetic term) can be derived from the closure structure, or whether it requires an additional ansatz --- and whether the WKB transport regime can be made precise enough to serve as an effective quantum mechanics for the closure framework.

%==================================================================
\bibliographystyle{plainnat}
\bibliography{references}

\end{document}
